\section{Apreciación de la herramienta}
\label{sec:apreciacion}

Se señalan dos características a favor y dos en contra de Flutter, cada una con la evidencia que dejó este proyecto.

\subsection{Dos características a favor}
\begin{enumerate}
	\item \textbf{Layout declarativo y componible, con un widget por decisión.} Alinear, repartir espacio, superponer, recortar o hacer scroll son widgets con nombre propio, y la regla única «las restricciones bajan, los tamaños suben» \cite{flutterconstr} permite razonar cualquier pantalla de la misma forma. En MatrLab esto se tradujo en reutilización real: \texttt{StatTile} se usa 8 veces en 3 pantallas, \texttt{ContenidoCentrado} 6 veces y \texttt{CupoBar} aparece tanto en el detalle del curso como en el del grupo, siempre con el mismo aspecto. Además, el catálogo de widgets de la sección~\ref{sec:catalogo} es pequeño (unos 18 widgets de layout) y cubrió las 9 pantallas.
	\item \textbf{Un solo código que se adapta, sin paquetes externos.} La misma aplicación pasa de teléfono a tableta cambiando únicamente la navegación, con un \texttt{LayoutBuilder} y un punto de corte de 600\,dp, y la cuadrícula de cursos se adapta sola con \texttt{SliverGridDelegateWithMaxCrossAxisExtent}. Todo el proyecto (2\,262 líneas en \texttt{lib/}) se hizo solo con el SDK: Material 3, tema, navegación, formularios, tablas y estado (\texttt{ChangeNotifier} + \texttt{InheritedNotifier}) vienen incluidos. Además, el desarrollo en Flutter admite recarga en caliente con estado conservado \cite{flutterarq}.
\end{enumerate}

\subsection{Dos características en contra}
\begin{enumerate}
	\item \textbf{Anidamiento profundo de widgets.} Como todo es un widget y cada ajuste (margen, alineación, recorte) es otro widget envolvente, el código se hunde hacia la derecha. En \texttt{cursos\_screen.dart} la tarjeta de curso llega a 18 niveles de sangría (\texttt{Card} $\rightarrow$ \texttt{InkWell} $\rightarrow$ \texttt{Padding} $\rightarrow$ \texttt{Row} $\rightarrow$ \texttt{Expanded} $\rightarrow$ \texttt{Column} $\rightarrow$ \texttt{Row}\ldots), lo que obliga a extraer clases y funciones para mantener la legibilidad; la propia documentación recomienda separar partes de la interfaz en variables o funciones \cite{flutterlayout}.
	\item \textbf{Las reglas de restricciones generan errores en tiempo de ejecución que no son evidentes.} Un \texttt{Expanded} dentro de un contenedor con scroll, un \texttt{GridView} dentro de un \texttt{ListView} o un \texttt{Card} sin tamaño producen desbordes, excepciones o widgets invisibles \cite{flutterconstr,flutterlayout}. En el panel hubo que recordar \texttt{shrinkWrap: true} y \texttt{NeverScrollableScrollPhysics} para evitar el error (tabla~\ref{tab:errores}). A esto se suma aprender Dart y su modelo de widgets, que para quien viene de JavaScript o C\# es una barrera inicial (puntaje 3 en la tabla~\ref{tab:decision}); y, a diferencia de React Native, Flutter no usa los controles del sistema sino los suyos \cite{flutterarq,rncomp}, de modo que la apariencia nativa depende del tema que se defina.
\end{enumerate}

\section{Conclusiones}
\label{sec:conclusiones}
\begin{itemize}
	\item Partir de un \textbf{árbol de problemas} ordenó el diseño: cada causa tiene una funcionalidad y cada funcionalidad una pantalla (figura~\ref{fig:causafun}), de modo que no hay pantallas sin justificación ni causas sin respuesta, salvo la concurrencia masiva (C3), que solo puede resolverse con un servidor.
	\item Entre los tres frameworks analizados, \textbf{Flutter} resultó el más adecuado para este laboratorio por su modelo de layout explícito y consistente (matriz de la tabla~\ref{tab:decision}); React Native queda como alternativa cercana si se prefiere JavaScript, y .NET MAUI si el equipo ya trabaja en el ecosistema .NET.
	\item El conocimiento clave de layout en Flutter cabe en una regla (restricciones abajo, tamaños arriba, posición decidida por el padre) y en un catálogo reducido de widgets; la mayoría de los errores se explican con esa regla.
	\item Se logró una aplicación completa sin paquetes externos, con lógica de negocio separada y probada, y con una interfaz que se adapta de teléfono a tableta.
	\item \textbf{Trabajo futuro:} conectar MatrLab a un servidor y una base de datos para resolver la concurrencia real (C3) y las notificaciones push; comparar el esfuerzo con el prototipo de Adalo del Laboratorio~2 \cite{lab2}; y probar la aplicación en dispositivos Android reales.
\end{itemize}
